\Needspace{8\baselineskip}
\Article{185}
The number of all reduced forms with a given determinant $D$ is always finite, and the forms themselves can be found in two ways. Let us denote all reduced forms with determinant $D$ in general by $(a,b,c)$, so that we have to determine all the values of $a,b,c$.

\textit{First method.} For $a$, take all numbers (both positive and negative) less than $2\surd D$ for which $D$ is a quadratic residue. For each $a$, set $b$ equal to all positive values of the expression $\surd D\ (\mathrm{mod.}\ a)$ lying between $\surd D$ and $\surd D\mp a$; and $c$ is to be set, for each pair of determined values of $a,b$, $=\frac{bb-D}{a}$. Any forms obtained this way in which $\pm a$ lies outside the interval between $\surd D+b$ and $\surd D-b$ must be discarded.

\textit{Second method.} For $b$, take all positive numbers less than $\surd D$. For each $b$, factor $bb-D$ in every possible way into two factors which, disregarding their signs, lie between $\surd D+b$ and $\surd D-b$, and set one $=a$ and the other $=c$. Each factorization clearly gives two forms, since each factor must be set both $=a$ and $=c$.
\par
\SourceLeaf{172}
\textit{Example.} Let $D=79$; then there are twenty-two values of $a$: ${\mp1}$, $2$, $3$, $5$, $6$, $7$, $9$, $10$, $13$, $14$, $15$. These give nineteen forms:
\[
\begin{gathered}
(1,8,-15),\quad (2,7,-15),\quad {(3,8,-5)},\quad (3,7,-10),\\
(5,8,-3),\quad (5,7,-6),\quad (6,7,-5),\quad (6,5,-9),\\
(7,4,-9),\quad (7,3,-10),\quad (9,5,-6),\quad (9,4,-7),\\
(10,7,-3),\quad (10,3,-7),\quad (13,1,-6),\quad (14,3,-5),\\
(15,8,-1),\quad (15,7,-2),\quad (15,2,-5)
\end{gathered}
\]
and just as many others obtained from them by changing the signs of the outer terms, namely $(-1,8,15)$, $(-2,7,15)$, etc., making thirty-eight in all. But six of these, $(\pm13,1,\mp6)$, $(\pm14,3,\mp5)$, $(\pm15,2,\mp5)$, must be discarded; the remaining thirty-two include all the reduced forms. The second method gives the same forms in the following order\HistoricalNote{*)}{For $b=1$, ${-78}$ cannot be factored into two factors which, disregarding their signs, lie between $\surd79+1$ and $\surd79-1$. So this value must be skipped, and for the same reason so must the values $2$ and $6$.}:
\[
\begin{gathered}
(\pm7,3,\mp10),\quad (\pm10,3,\mp7),\quad {\pm7,4,\mp9)},\quad (\pm9,4,\mp7),\\
(\pm6,5,\mp9),\quad (\pm9,5,\mp6),\quad (\pm2,7,\mp15),\quad (\pm3,7,\mp10),\\
(\pm5,7,\mp6),\quad (\pm6,7,\mp5),\quad (\pm10,7,\mp3),\quad (\pm15,7,\mp2),\\
(\pm1,8,\mp15),\quad (\pm3,8,\mp5),\quad (\pm5,8,\mp3),\quad (\pm15,8,\mp1).
\end{gathered}
\]

\Needspace{8\baselineskip}
\Article{186}
Let $F$ be a reduced form with determinant $D$, and let $F'$ be its reduced right neighbor; let $F''$ in turn be the reduced right neighbor of this form, $F'''$ the reduced right neighbor of $F''$, and so on. All the forms $F',F'',F'''$ etc. are then completely determined, and are properly equivalent both to one another and to $F$. But since the number of all reduced forms with a given determinant is finite, the forms in the infinite sequence $F,F',F''$ etc. cannot all be different. Suppose $F^m$ and $F^{m+n}$ are identical. Then $F^{m-1}$, $F^{m+n-1}$ are reduced left neighbors of the same reduced form, and so are identical; in the same way $F^{m-2}$ and $F^{m+n-2}$ etc., and finally $F$ and $F^n$, are identical. So if the sequence $F,F',F''$ etc. is continued far enough, the first form $F$ must eventually recur. If we suppose that $F^n$ is the first form identical to $F$, or that all $F',F''\mathinner{\ldotp\ldotp\ldotp\ldotp} F^{n-1}$ are different from $F$, we easily see that all the forms $F,F',F''\mathinner{\ldotp\ldotp\ldotp\ldotp} F^{n-1}$ are different. We will call this collection of \SourceLeaf{173}forms the \textit{period of the form} $F$. If the sequence is continued beyond the last form of the period, the same forms $F,F',F''$ etc. appear again, and the whole infinite sequence $F,F',F''$ etc. consists of this period of $F$ repeated infinitely often.

The sequence $F,F',F''$ etc. can also be continued backwards by placing before $F$ the reduced form ${}'F$, its left neighbor; before this, the reduced form ${}''F$, its left neighbor, and so on. This gives a sequence of forms infinite \textit{in both directions}:
\[
\mathinner{\ldotp\ldotp\ldotp\ldotp} {}'''F,\ {}''F,\ {}'F,\ F,\ F',\ F'',\ F'''\mathinner{\ldotp\ldotp\ldotp\ldotp}
\]
We easily see that ${}'F$ is identical to $F^{n-1}$, ${}''F$ to $F^{n-2}$, etc., so the sequence on the left also consists of the period of $F$ repeated infinitely often.

If we assign to the forms $F,F',F''$ etc., ${}'F$, ${}''F$ etc. the indices $0$, $1$, $2$ etc., $-1$, $-2$ etc., and in general to $F^m$ the index $m$ and to ${}^mF$ the index $-m$, then clearly \textit{any two forms of the sequence are identical or different according as their indices are congruent or incongruent modulo $n$.}

\textit{Example.} The period of the form $(3,8,-5)$, with determinant $=79$, is found to be ${3,8,-5)}$, $(-5,7,6)$, $(6,5,-9)$, $(-9,4,7)$, $(7,3,-10)$, $(-10,7,3)$. After the last, $(3,8,-5)$ appears again. So here $n=6$.

\Needspace{8\baselineskip}
\Article{187}
Here are some general observations about these periods.

1) If the forms $F,F',F''$ etc.; ${}'F,{}''F,{}'''F$ etc. are written as
\[
\begin{gathered}
{a,b,-a')},\quad (-a',b',a''),\quad (a'',b'',-a''')\text{ etc.};\\
(-{}'a,{}'b,a),\quad ({}''a,{}''b,-{}'a),\quad (-{}'''a,{}'''b,{}''a)
\end{gathered}
\]
then all $a,a',a'',a'''$ etc., ${}'a,{}''a,{}'''a$ etc. have \textit{the same sign} (Article~184,~1), while all $b,b',b''$ etc., ${}'b,{}''b$ etc. are positive.

2) This shows that the number $n$ (the number of forms making up the period of $F$) is always \textit{even}. For the first term of any form $F^m$ in this period clearly has the same sign as the first term $a$ of $F$ if $m$ is even, and the opposite sign if $m$ is odd. Since $F^n$ and $F$ are identical, $n$ must be even.
\SourceLeaf{174}

\Needspace{14\baselineskip}
3) By Article~184,~6, the algorithm for finding the numbers $b',b'',b'''$ etc., $a'',a'''$ etc. is as follows:
\[
\begin{array}{lll}
b'\equiv-b\ (\mathrm{mod.}\ a')&\text{between }\surd D\text{ and }\surd D\mp a';&a''=\dfrac{D-b'b'}{a'}\\[6pt]
b''\equiv-b'\ (\mathrm{mod.}\ a'')&\mathinner{\ldotp\ldotp\ldotp\ldotp\ldotp\ldotp\ldotp\ldotp\ldotp}\quad\surd D\mp a'';&a'''=\dfrac{D-b''b''}{a''}\\[6pt]
b'''\equiv-b''\ (\mathrm{mod.}\ a''')&\mathinner{\ldotp\ldotp\ldotp\ldotp\ldotp\ldotp\ldotp\ldotp\ldotp}\quad\surd D\mp a''';&a''''=\dfrac{D-b'''b'''}{a'''}
\end{array}
\]
\centerline{etc.}
In the second column we take the upper or lower signs according as $a,a',a''$ etc. are positive or negative. Instead of the formulas in the third column, we may also use the following, which are more convenient when $D$ is large:
\[
\begin{aligned}
a''&=\frac{b+b'}{a'}(b-b')+a\\
a'''&=\frac{b'+b''}{a''}(b'-b'')+a'\\
a''''&=\frac{b''+b'''}{a'''}({b''}-b''')+a''
\end{aligned}
\]
\centerline{etc.}

4) Any form $F^m$ contained in the period of $F$ has essentially the same period as $F$. Its period is $F^m$, $F^{m+1}\mathinner{\ldotp\ldotp\ldotp\ldotp} F^{n-1}$, $F,F'\mathinner{\ldotp\ldotp\ldotp\ldotp} F^{m-1}$, containing the same forms in the same order as the period of $F$, and differing from it only in where it begins and ends.

5) So all reduced forms with the same determinant $D$ can be \textit{distributed} into periods. Take any one of these forms, $F$, and find its period $F,F',F''\mathinner{\ldotp\ldotp\ldotp\ldotp} F^{n-1}$, which we will denote by $P$. If this does not yet include all reduced forms with determinant $D$, let $G$ be one not contained in it, and $Q$ its period. Then $P$ and $Q$ clearly cannot have a form in common; otherwise $G$ would also have to be contained in $P$, and the periods would coincide entirely. If $P$ and $Q$ do not yet exhaust all reduced forms, one of the missing forms, $H$, supplies a third period $R$, having no form in common with either $P$ or $Q$. We can continue this way until all the reduced forms are exhausted. For \textit{example}, all reduced forms with determinant $79$ are distributed into six periods:
\SourceLeaf{175}
\[
\begin{array}{rl}
\mathrm{I.}&(1,8,-15),\quad (-15,7,2),\quad (2,7,-15),\quad (-15,8,1).\\
\mathrm{II.}&(-1,8,15),\quad (15,7,-2),\quad (-2,7,15),\quad (15,8,-1).\\
\mathrm{III.}&(3,8,-5),\quad (-5,7,6),\quad (6,5,-9),\\
&(-9,4,7),\quad (7,3,-10),\quad (-10,7,3).\\
\mathrm{IV.}&(-3,8,5),\quad (5,7,-6),\quad (-6,5,9),\\
&(9,4,-7),\quad (-7,3,10),\quad (10,7,-3).\\
\mathrm{V.}&(5,8,-3),\quad (-3,7,10),\quad (10,3,-7),\\
&(-7,4,9),\quad (9,5,-6),\quad (-6,7,5).\\
\mathrm{VI.}&(-5,8,3),\quad (3,7,-10),\quad (-10,3,7),\\
&(7,4,-9),\quad (-9,5,6),\quad (6,7,-5).
\end{array}
\]

6) Let us call forms made up of the same terms in reverse order, such as $(a,b,-a')$, $(-a',b,a)$, \textit{associated forms}. By Article~184,~7 we easily see that if the period of the reduced form $F$ is $F,F',F''\mathinner{\ldotp\ldotp\ldotp\ldotp} F^{n-1}$, if $F$ has the associated form $f$, and if the forms $F^{n-1},F^{n-2}\mathinner{\ldotp\ldotp\ldotp\ldotp} F'',F'$ have the respective associated forms $f',f''\mathinner{\ldotp\ldotp\ldotp\ldotp} f^{n-2},f^{n-1}$, then the period of $f$ is $f,f',f''\mathinner{\ldotp\ldotp\ldotp\ldotp} f^{n-2},f^{n-1}$. So it consists of just as many forms as the period of $F$. We will call the periods of associated forms \textit{associated periods}. In our example, periods III and VI are associated, as are IV and V.

7) But the form $f$ may itself occur in the period of its associated form $F$, as in periods I and II of our example. Then the period of $F$ agrees with that of $f$, or \textit{the period of $F$ is associated with itself}. Whenever this happens, two ambiguous forms occur in this period. Suppose the period of $F$ consists of $2n$ forms, or that $F$ and $F^{2n}$ are identical; and let $2m+1$ be the index of $f$ in the period of $F$\HistoricalNote{*)}{This index must be odd, since the first terms of the forms ${F},f$ clearly have opposite signs (\textit{see above}, $2$).}, so that $F^{2m+1}$ and $F$ are associated. Then $F'$ and $F^{2m}$ are also associated, as are $F''$ and $F^{2m-1}$ etc., and so also $F^m$ and $F^{m+1}$. Let $F^m=(a^m,b^m,-a^{m+1})$, $F^{m+1}=(-a^{m+1},b^{m+1},a^{m+2})$. Then $b^m+b^{m+1}\equiv0\ (\mathrm{mod.}\ a^{m+1})$; but by the definition of associated forms, $b^m=b^{m+1}$, and so $2b^{m+1}\equiv0\ (\mathrm{mod}\ a^{m+1})$, or the form $F^{m+1}$ is ambiguous. --- In the same way $F^{2m+1}$ and $F^{2n}$ are associated; so are $F^{2m+2}$ and $F^{2n-1}$, $F^{2m+3}$ and $F^{2n-2}$ etc., and finally $F^{m+n}$ and $F^{m+n+1}$. The latter is ambiguous, as a similar argument easily proves. But $m+1$ and $m+n+1$ are incongruent modulo $2n$, so the forms $F^{m+1}$ and $F^{m+n+1}$ are not identical (Article~186, where $n$ denotes what $2n$ denotes here). Thus in I the ambiguous forms are $(1,8,-15)$, $(2,7,-15)$, and in II they are $(-1,8,15)$, $(-2,7,15)$.
\SourceLeaf{176}

8) Conversely, \textit{every period containing an ambiguous form is associated with itself.} For if $F^m$ is a reduced ambiguous form, we easily see that its associated form (which is also reduced) is also its left neighbor, \textit{that is}, $F^{m-1}$ and $F^m$ are associated. But then the whole period is associated with itself. --- This shows that \textit{a period cannot contain exactly one ambiguous form.}

9) But \textit{it cannot contain more than two either.} Suppose the period of $F$, consisting of $2n$ forms, contains three ambiguous forms $F^\lambda,F^\mu,F^\nu$, with respective indices $\lambda,\mu,\nu$, so that $\lambda,\mu,\nu$ are distinct numbers lying between $0$ and $2n-1$ inclusive. Then the forms $F^{\lambda-1}$ and $F^\lambda$ are associated; so are $F^{\lambda-2}$ and $F^{\lambda+1}$ etc., and finally $F$ and $F^{2\lambda-1}$. For the same reason $F$ and $F^{2\mu-1}$ are associated, as are $F$ and $F^{2\nu-1}$. So $F^{2\lambda-1}$, $F^{2\mu-1}$, $F^{2\nu-1}$ are identical, and the indices $2\lambda-1$, $2\mu-1$, $2\nu-1$ are congruent modulo $2n$, and consequently $\lambda\equiv\mu\equiv\nu\ (\mathrm{mod.}\ n)$. \textit{This is absurd}, since between $0$ and $2n-1$ there clearly cannot lie three distinct numbers congruent modulo $n$.
